\chapter{Experiments}
\label{chapter:experiments}

\section{Research questions and hypotheses}
\label{section:hypotheses}

\begin{description}[leftmargin=1.2cm, labelwidth=1cm]
  \item[RQ1] \emph{Quality.} Does the hybrid mode reach at least the accuracy and macro-F1 of the LLM-only mode, and
        exceed the deterministic mode, on applicable requirements?
  \item[RQ2] \emph{Safety.} Which mode minimises the false-compliance rate, and how much human workload (needs-review
        rate) does that cost?
  \item[RQ3] \emph{Cost.} How many LLM calls, tokens, dollars and seconds does each mode need per audit, and how much does
        the cache save on a re-audit?
  \item[RQ4] \emph{Safeguards.} What do the Verifier, the LLM support check and the model size contribute (ablations)?
\end{description}
\noindent The \emph{independent variables} are the mode, the model, and the safeguard switches; the \emph{dependent
variables} are the metrics of Section~\ref{section:formal}.

\section{Data: a hand-labelled benchmark}
\label{section:data}

There is no public dataset of application repositories labelled against the AI Act, so we built one
(\texttt{app/benchmark}). Five small but realistic repositories cover the four risk tiers and different stacks
(Table~\ref{tab:apps}); each contains code, configuration, infrastructure-as-code, documentation, model cards and sample
logs as real teams would write them. Every one of the 45 curated requirements was labelled for every app (225 labels,
165 applicable), with a justification and file references. Labels were written \emph{before} running the rule engine, by
reading the repository against the requirement statement; claims without backing (README marketing, TODOs, commented code)
are not evidence; \stat{na} is used only when the requirement's risk tiers exclude the app's tier, so it is never a
judgement call. Table~\ref{tab:labels} gives the label distribution.

\begin{table}[htbp]
\caption{Benchmark applications.}
\label{tab:apps}
\centering\small
\begin{tabularx}{\linewidth}{@{}lllX@{}}
\toprule
\textbf{App} & \textbf{Domain} & \textbf{AI Act tier} & \textbf{Compliance profile} \\
\midrule
\texttt{hireright-screening} & CV ranking & high (Annex III 4(a)) & mostly compliant; logging via a custom journal, kill switch with an unusual name; gaps: bias evaluation, ROPA, DPIA \\
\texttt{quickhire-ranker} & CV ranking MVP & high (Annex III 4(a)) & mostly non-compliant: auto-rejection, protected attributes as features, no auth; decoys (commented logger, ``fully GDPR compliant'' README) \\
\texttt{supportbot-genai} & support chatbot & limited (Art.~50) & disclosure banner, erasure endpoint; gaps: unlabelled AI replies (TODO only), 2-year transcript retention \\
\texttt{credit-lens} & credit scoring & high (Annex III 5(b)) & strong governance; gaps: 30-day log retention, automatic declines up to EUR 25\,000, no instructions for use \\
\texttt{spamshield} & spam filter & minimal & most AI Act duties N/A; GDPR largely compliant; gaps: incident procedure, on-request erasure \\
\bottomrule
\end{tabularx}
\end{table}

\begin{table}[htbp]
\caption{Ground-truth label distribution (45 requirements $\times$ 5 apps).}
\label{tab:labels}
\centering\small
\begin{tabular}{@{}llrrrrr@{}}
\toprule
\textbf{App} & \textbf{Tier} & \stat{c} & \stat{p} & \stat{nc} & \stat{na} & \textbf{total} \\
\midrule
credit-lens & high & 32 & 7 & 3 & 3 & 45 \\
hireright-screening & high & 37 & 1 & 4 & 3 & 45 \\
quickhire-ranker & high & 1 & 2 & 39 & 3 & 45 \\
spamshield & minimal & 15 & 2 & 1 & 27 & 45 \\
supportbot-genai & limited & 15 & 2 & 4 & 24 & 45 \\
\midrule
\textbf{all} & & \textbf{100} & \textbf{14} & \textbf{51} & \textbf{60} & \textbf{225} \\
\bottomrule
\end{tabular}
\end{table}

The benchmark is deliberately adversarial for rules: \emph{indirect evidence} (regex false negatives, e.g.\ a
\texttt{DecisionJournal} instead of \texttt{logger}, ``pseudonymous'' instead of ``pseudonymise''), \emph{decoys} (regex
false positives, e.g.\ \texttt{\# logger.info(...)}, \texttt{override\_allowed = False}, a TODO to label AI replies) and
three requirements without any check. It is also adversarial for LLMs: plausible but unbacked claims in READMEs and
disproportionate-but-present controls (retention defined but too long) that require reading thresholds.

\subsection{Exploratory data analysis}
\label{section:eda}

Figures~\ref{fig:eda-corpus}--\ref{fig:eda-signals} describe the inputs. The AI Act dominates the corpus, with long and
structurally varied clauses (Annex~III, Art.~9, Art.~13). Most requirements have exactly one deterministic check, and three
have none (Figure~\ref{fig:eda-corpus}). The applications differ in the mix of evidence kinds: \texttt{hireright} and
\texttt{credit-lens} are documentation-heavy, \texttt{quickhire} is mostly code (Figure~\ref{fig:eda-bench}). The labels
are imbalanced per application by design (Figure~\ref{fig:eda-labels}). The signal heat map
(Figure~\ref{fig:eda-signals}) shows why a signal alone is weak evidence: for example, 50\% of \texttt{quickhire} chunks
carry a PII signal, yet that application has no data protection.

\begin{figure}[htbp]
\centering
\maybefig[0.49\linewidth]{eda_clauses}{Clauses per article and clause length.}\hfill
\maybefig[0.49\linewidth]{eda_requirements}{Requirements per category, severity, checks per requirement.}
\caption{Policy corpus: clauses per article and clause length (left); requirements per category, severity and number of
deterministic checks (right).}
\label{fig:eda-corpus}
\end{figure}

\begin{figure}[htbp]
\centering
\maybefig[0.8\linewidth]{eda_benchmark_files}{Files and evidence chunks per kind and app.}
\caption{Benchmark repositories: files and evidence chunks per evidence kind.}
\label{fig:eda-bench}
\end{figure}

\begin{figure}[htbp]
\centering
\maybefig[0.8\linewidth]{eda_labels}{Labels per app and applicable labels per category.}
\caption{Ground-truth labels per application (left) and applicable labels per category (right).}
\label{fig:eda-labels}
\end{figure}

\begin{figure}[htbp]
\centering
\maybefig[0.6\linewidth]{eda_signals}{Share of evidence chunks carrying each signal, per app.}
\caption{Share of evidence chunks carrying each collector signal, per application.}
\label{fig:eda-signals}
\end{figure}

\section{Protocol}
\label{section:protocol}

\begin{itemize}[leftmargin=*]
  \item \textbf{Configurations.} \emph{deterministic}; \emph{hybrid} and \emph{llm} with Claude Sonnet (default) and
        Claude Haiku (smaller, cheaper); all three policies at once on each app (45 requirements per audit).
  \item \textbf{Runs.} \texttt{app/benchmark/evaluate.py} audits the five apps in parallel processes (3 concurrent LLM
        calls per audit), stores every full report, the per-item predictions and a $5\times5$ confusion matrix per
        configuration. Temperature cannot be set through the CLI; variance is estimated by repeating cold (uncached)
        runs. \resultsTODO{number of repetitions and seeds.}
  \item \textbf{Ablations.} (a) Verifier disabled (\texttt{--no-verifier}); (b) LLM claim-support check enabled
        (\texttt{--verify-llm}); (c) model size (Sonnet vs.\ Haiku); (d) routing threshold $\tau$ / routing of confident
        \stat{nc} rule verdicts; (e) cache: cold run vs.\ warm re-audit of unchanged code.
  \item \textbf{Statistics.} 95\% bootstrap confidence intervals over items (1\,000 resamples) for accuracy, macro-F1 and
        FCR, and McNemar's exact test on paired per-item correctness between modes.
  \item \textbf{Hardware.} \todo{CPU / RAM / OS of the machine}; LLM inference runs on Anthropic's servers via the
        Claude Code CLI \todo{CLI version}, so local memory is dominated by the Python process.
\end{itemize}

\section{Results}
\label{section:results}

\IfFileExists{results_tables.tex}{% Main results. Source: app/benchmark/results/SUMMARY.md (v0_baseline.json, v1_improved.json) and
% app/benchmark/results/stats_tests.json (app/benchmark/stats_tests.py).
Table~\ref{tab:main-results} gives all configurations on the 165 applicable labels; Table~\ref{tab:stats} adds 95\%
bootstrap confidence intervals and paired McNemar tests.

\begin{table}[htbp]
\caption{Results on the 165 applicable labels (percentages). Cost is the equivalent USD reported by the CLI for the
five audits (cached calls counted at their original cost); s/app is the wall-clock per app audit of a cold run.
v0 = baseline, v1 = improved system (Section~\ref{section:v1}). Best v1 value per column in bold.}
\label{tab:main-results}
\centering\scriptsize
\setlength{\tabcolsep}{3.2pt}
\resizebox{\linewidth}{!}{%
\begin{tabular}{@{}llrrrrrrrrrr@{}}
\toprule
& \textbf{Config} & \textbf{Acc} & \textbf{M-F1} & \textbf{Gap P} & \textbf{Gap R} & \textbf{FCR}$\downarrow$ & \textbf{NR} & \textbf{Acc@def} & \textbf{Calls} & \textbf{Cost \$} & \textbf{s/app} \\
\midrule
\multirow{3}{*}{v0} & deterministic & 79.4 & 72.3 & 78.1 & 87.7 & 12.3 & 5.5 & 84.0 & 0 & 0.00 & 0.2 \\
& hybrid (Sonnet) & 86.1 & 77.1 & 83.8 & 87.7 & 12.3 & 0.0 & 86.1 & 23 & 0.67 & 25 \\
& llm (Sonnet) & 71.5 & 67.6 & 60.8 & \textbf{100.0} & \textbf{0.0} & 0.6 & 72.0 & 170 & 4.81 & 94 \\
\midrule
\multirow{5}{*}{v1} & deterministic & 82.4 & 74.9 & 80.3 & 93.8 & 6.2 & 5.5 & 87.2 & 0 & \textbf{0.00} & \textbf{0.2} \\
& \textbf{hybrid (Sonnet)}, $\tau=0.675$ & \textbf{94.5} & \textbf{86.6} & \textbf{96.9} & 96.9 & 3.1 & 0.6 & \textbf{95.1} & 103 & 2.77 & 50 \\
& llm (Sonnet) & 90.9 & 83.1 & 85.5 & \textbf{100.0} & \textbf{0.0} & 0.6 & 91.5 & 170 & 4.68 & 91 \\
& hybrid (Haiku) & 83.6 & 75.4 & 76.8 & 96.9 & 3.1 & 5.5 & 88.5 & 103 & 0.19 & 75 \\
& llm (Haiku) & 71.5 & 66.7 & 61.9 & \textbf{100.0} & \textbf{0.0} & 7.3 & 77.1 & 170 & 0.32 & 116 \\
\midrule
\multirow{3}{*}{\shortstack[l]{v1\\ablation}} & hybrid (Sonnet), route all ($\tau>1$) & 90.9 & 82.1 & 87.7 & 98.5 & 1.5 & 0.6 & 91.5 & 170 & 4.82 & 90 \\
& hybrid (Sonnet), no Verifier & 94.5 & 86.6 & 96.9 & 96.9 & 3.1 & 0.6 & 95.1 & 103 & 2.77 & 50 \\
& hybrid (Sonnet) + LLM claim verifier & 88.5 & 83.1 & 85.1 & 96.9 & 3.1 & 6.7 & 94.8 & 156 & 3.89 & 89 \\
\bottomrule
\end{tabular}}
\end{table}

\noindent In every configuration the citation validity was 100\% (no cited id outside the evidence index) and in every v1
configuration that uses the LLM, every \stat{c} verdict cites at least one evidence item. Risk-tier
accuracy is 100\% for all v1 configurations (v0 deterministic: 80\%, \texttt{spamshield} classified as limited).

\begin{table}[htbp]
\caption{Statistical validation: 95\% percentile bootstrap CIs (10\,000 resamples of the 165 items) and exact
two-sided McNemar tests on paired per-item correctness ($b$: only A correct, $c$: only B correct).}
\label{tab:stats}
\centering\scriptsize
\begin{tabular}{@{}lccc@{}}
\toprule
\textbf{Configuration} & \textbf{Accuracy} & \textbf{Macro-F1} & \textbf{FCR} \\
\midrule
v0 deterministic & 79.4 [72.7, 85.5] & 72.3 [61.4, 81.6] & 12.3 [4.8, 20.9] \\
v0 llm (Sonnet) & 71.5 [64.8, 78.2] & 67.6 [61.0, 73.8] & 0.0 [0.0, 0.0] \\
v1 deterministic & 82.4 [76.4, 87.9] & 74.9 [64.2, 84.0] & 6.2 [1.4, 12.7] \\
v1 hybrid (Sonnet) & 94.5 [90.9, 97.6] & 86.6 [77.0, 94.1] & 3.1 [0.0, 7.9] \\
v1 llm (Sonnet) & 90.9 [86.7, 95.2] & 83.1 [74.6, 90.1] & 0.0 [0.0, 0.0] \\
v1 hybrid (Haiku) & 83.6 [77.6, 89.1] & 75.4 [66.8, 83.3] & 3.1 [0.0, 7.9] \\
\bottomrule
\end{tabular}
\vspace{0.6em}

\begin{tabular}{@{}llrrr@{}}
\toprule
\textbf{A} & \textbf{B} & $b$ & $c$ & \textbf{$p$-value} \\
\midrule
v1 hybrid (Sonnet) & v1 deterministic & 22 & 2 & $3.6\times10^{-5}$ \\
v1 hybrid (Sonnet) & v1 llm (Sonnet) & 11 & 5 & 0.21 \\
v1 hybrid (Sonnet) & v1 hybrid, route all & 7 & 1 & 0.070 \\
v1 hybrid (Sonnet) & v1 hybrid (Haiku) & 20 & 2 & $1.2\times10^{-4}$ \\
v1 hybrid (Sonnet) & v1 hybrid + LLM claim verifier & 10 & 0 & 0.0020 \\
v1 llm (Sonnet) & v0 llm (Sonnet) & 34 & 2 & $1.9\times10^{-8}$ \\
v1 hybrid (Sonnet) & v0 hybrid (Sonnet) & 16 & 2 & 0.0013 \\
v1 deterministic & v0 deterministic & 5 & 0 & 0.063 \\
\bottomrule
\end{tabular}
\end{table}
}{%
Table~\ref{tab:main-results} is filled from \texttt{app/benchmark/results/SUMMARY.md}.

\begin{table}[htbp]
\caption{Main results on the 165 applicable labels (percentages; cost and latency per app audit).}
\label{tab:main-results}
\centering\scriptsize
\setlength{\tabcolsep}{3.5pt}
\begin{tabular}{@{}lrrrrrrrrrr@{}}
\toprule
\textbf{Config} & \textbf{Acc} & \textbf{M-F1} & \textbf{Gap P} & \textbf{Gap R} & \textbf{FCR}$\downarrow$ & \textbf{NR} & \textbf{Acc@def} & \textbf{Cit.\ valid} & \textbf{LLM calls} & \textbf{Cost \$ / s} \\
\midrule
deterministic        & \res{} & \res{} & \res{} & \res{} & \res{} & \res{} & \res{} & \res{} & 0 & 0 / \res{} \\
hybrid (Sonnet)      & \res{} & \res{} & \res{} & \res{} & \res{} & \res{} & \res{} & \res{} & \res{} & \res{} \\
llm (Sonnet)         & \res{} & \res{} & \res{} & \res{} & \res{} & \res{} & \res{} & \res{} & \res{} & \res{} \\
hybrid (Haiku)       & \res{} & \res{} & \res{} & \res{} & \res{} & \res{} & \res{} & \res{} & \res{} & \res{} \\
\midrule
hybrid, no Verifier  & \res{} & \res{} & \res{} & \res{} & \res{} & \res{} & \res{} & \res{} & \res{} & \res{} \\
hybrid, LLM verifier & \res{} & \res{} & \res{} & \res{} & \res{} & \res{} & \res{} & \res{} & \res{} & \res{} \\
\bottomrule
\end{tabular}
\end{table}

\resultsTODO{fill Table~\ref{tab:main-results} (or provide \texttt{docs/report/results\_tables.tex}, which replaces this
block automatically), add bootstrap CIs and McNemar p-values.}
}

\begin{figure}[htbp]
\centering
\maybefig{mode_comparison}{Bar chart: accuracy, macro-F1, gap F1, FCR, needs-review per configuration.}
\caption{Quality metrics per configuration.}
\label{fig:mode-comparison}
\end{figure}

\begin{figure}[htbp]
\centering
\maybefig[0.48\linewidth]{confusion_hybrid}{Confusion matrix (truth $\times$ prediction) of the hybrid mode.}\hfill
\maybefig[0.48\linewidth]{confusion_llm}{Confusion matrix of the LLM-only mode.}
\caption{Confusion matrices (rows: ground truth, columns: prediction) for hybrid (left) and LLM-only (right).}
\label{fig:confusion}
\end{figure}

\begin{figure}[htbp]
\centering
\maybefig[0.7\linewidth]{cost_quality}{Scatter: cost (USD) or LLM calls per audit vs.\ accuracy / FCR, one point per configuration.}
\caption{Cost--quality trade-off of the configurations.}
\label{fig:cost-quality}
\end{figure}

\begin{figure}[htbp]
\centering
\maybefig{per_category}{Heat map or grouped bars: accuracy per requirement category and mode.}
\caption{Accuracy per requirement category.}
\label{fig:per-category}
\end{figure}

\section{Complexity and sustainability}
\label{section:complexity}

\paragraph{Time and memory.} The deterministic mode is CPU-only and runs in well under a second per application; memory
is bounded by the evidence index (a few thousand chunks for our apps). In LLM modes wall-clock time is dominated by the
judge calls and scales with the number of routed requirements divided by the concurrency.
\resultsTODO{per-stage timings (requirements, evidence, risk, deterministic, LLM judge, report) and peak RSS.}

\paragraph{Cost and energy.} Cost is measured as the equivalent USD reported by the CLI per call
(\texttt{total\_cost\_usd}; the subscription itself is a flat fee) together with input/output tokens. We cannot meter the
energy of a hosted model directly, so tokens and the number of calls are used as proxies: inference energy grows with
generated and processed tokens~\cite{luccioni2024power,strubell2019energy}. Three design choices reduce them: (1)
\emph{routing} -- the LLM sees only the requirements the rules cannot settle; (2) \emph{bounded prompts} -- at most 14
evidence items and 12\,000 characters per judgement, one requirement per call, no tool use; (3) a \emph{content-addressed
cache}, so a re-audit of an unchanged repository (the common CI case) costs no new calls.
\resultsTODO{cold vs.\ warm cost and calls per audit; tokens per configuration.}

\section{Explainability}
\label{section:explainability}

Every verdict is explainable by construction: it carries the rule-check trace (which check passed or failed and where),
the judge's rationale that refers to evidence aliases, the cited \texttt{file:line} snippets, the gaps and the
remediation; conflicts keep both the rule and the LLM rationale. \texttt{ccauditor explain} and the dashboard side sheet
render this justification for a single requirement, and SARIF places the finding on the cited lines in code review.

\section{Discussion}
\label{section:discussion}

\resultsTODO{answer RQ1--RQ4 with the numbers: is the hybrid hypothesis supported; where does each mode fail (per
category, per app); effect of the Verifier and of model size; statistical significance; cost per correct verdict.}

Independently of the final numbers, the design implies the following expected failure modes, which the analysis should
confirm or refute: (i) the rule engine is precise on decoy-free, syntactic requirements but propagates confident false
negatives on indirect evidence (Section~\ref{section:example}); (ii) the LLM-only mode recovers indirect evidence but
tends to be stricter, flagging more items as gaps (higher gap recall, lower gap precision) and costing an order of
magnitude more calls; (iii) the hybrid mode inherits the rule engine's confident errors because they are never routed,
which bounds its accuracy from above and is the main target for threshold tuning; (iv) data-governance and
human-oversight requirements are the hardest for rules, because their evidence is mostly prose.

\paragraph{Threats to validity.} The benchmark is small (5 apps, 225 labels), synthetic and labelled by the authors who
also wrote the rules, which risks optimistic results for the deterministic mode on the patterns we anticipated; we
mitigated this by labelling before running the rules and by planting traps against our own rules. Policies are selected
provisions, requirements are curated, and the AI Act will be complemented by harmonised standards that may change what
counts as sufficient evidence.
